\documentclass[a4paper]{article}
\usepackage{lab}
\usepackage{subcaption}
\usepackage{tabularx}
\usepackage{multirow}

\begin{document}
\begin{titlepage}
    \thispagestyle{fancy} % Makes the page fancy
    \rhead{}
    \chead{}
    \lhead{ECGR 3120-001} % For the left header
    \rfoot{}
    \cfoot{}
    \lfoot{}
    \renewcommand{\headrulewidth}{0pt}
    \begin{center}
        \vspace*{2cm}


        \Large\textbf{University of North Carolina at Charlotte \\ [1pt]
        Department of Electrical and Computer Engineering \\ [2pt]
        \emph{ECGR 3120 Inkjet Printer Simulation} \\ [7pt]
        }

        \noindent\rule{15cm}{0.4pt}
        \vfill
        \vspace{3cm}
        \textit{Group Project Report} \\
        Team Members: Govind Menon,\\ Adi Chettri, and\\ Rushil Dasari \\
        Date: December 5, 2025 \\

        \vspace{1cm}

        \vfill
        \vfill
        \vfill

    \end{center}
\end{titlepage}

\section*{Q1}
\subsection*{Simulate the motion of the droplet when no voltage is applied across the capacitor. How much time would it take for the droplet to reach the center of the paper?}\par

The droplet is travelling along the x-axis towards the paper. Because there is no acceleration in the x-direction: \\ \\
$V_X = 20m/s\ (constant)$ which is independent of $V_Y$.\\ \\
The path taken by each droplet regardless of the voltage applied is the same in the x direction. This displacement can be calculated by integrating $V_X$ which gives us:\\ \\
$S_X = 20*t\ m$. Solving for time gives us: $\frac{S_X}{20} = t\ sec$. It is given that the distance from droplet gun to the paper is 3mm.\par
Therefore, the time taken by a droplet to reach the center of the paper with no voltage applied across the capacitor is \textbf{0.15 milliseconds}.

\section*{Q2}
\subsection*{Simulate the act of drawing a vertical line (similar to the letter “I” ) along the center of the paper at 300 dpi resolution. How much time would it take to draw the letter “I”?}\par


\end{document}